\Subhead{4.6.2}{TE-, TH- и TEM-волны. Критические частота и длина волновода. Диэлектрические волноводы. Периодические структуры и замедляющие системы}

В полом волноводе поперечная собственная задача даёт $\beta^2=\omega^2\varepsilon\mu-k_c^2$; критические частота и длина волны $\omega_c=k_c/\sqrt{\varepsilon\mu}$, $\lambda_c=2\pi/k_c$ в заполняющей среде. Под «критической длиной» в формулировке программы здесь понимается критическая длина волны $\lambda_c$, а не длина отрезка волновода. В программе написано «TH»; здесь используется общепринятое TM для мод с $H_z=0$. TE имеет $E_z=0$, TEM — обе продольные компоненты нулевые и требует многосвязной структуры, например коаксиала. Для TE/TM волновые сопротивления различны и зависят от $\beta$; это не всегда отношение полного напряжения к току, как у TEM-линии. Диэлектрические световоды допускают гибридные моды; их основные условия разобраны в вопросе~\Qref{2.6}.

Обычная волноводная фазовая скорость больше скорости в веществе, поэтому для взаимодействия с нерелятивистским электронным пучком нужна \emph{замедляющая система}. В спирали поле движется вдоль длинного проводящего пути, а его осевой фазовый набег становится большим. Периодические гребёнки и связанные резонаторы дают блоховские волны
\begin{equation}
 E_z(z)=\sum_m E_m\e^{\ii(\beta+2\pi m/d)z},
\end{equation}
где $d$ — период. Отдельная пространственная гармоника синхронна с пучком при $\omega=(\beta+2\pi m/d)v_0$. Фазовая и групповая скорости такой структуры различаются по величине и иногда по знаку; это используется для ЛБВ и ЛОВ. Число ячеек, дисперсия, сопротивление связи и потери определяют эффективность взаимодействия.
